% Level 2, Week 9 -- Event-Driven Games (student workbook)
\chapterweek{9}{Event-Driven Games}{Keys and timers make things happen; a game loop moves the world one frame at a time.}

\section{What you will learn}
\begin{itemize}[leftmargin=1.2em]
  \item How to make keys call your functions with \texttt{onkey} and \texttt{listen}.
  \item How a \textbf{game loop} with \texttt{ontimer} moves everything one frame at a time.
  \item How to keep the game's \textbf{state} in a dictionary and check if two things are \textbf{touching}.
  \item How to start planning your own final project.
\end{itemize}

\section{The big idea: games wait for events}
Most programs you wrote ran from top to bottom and stopped. A game \textbf{waits}. When something happens (a key press, a timer ringing),
turtle calls one of your functions. The thing that happens is an \textbf{event}; the function that answers it is a \textbf{handler}.

\begin{pycode}[Keys call handlers]
def go_left():
    player.setx(player.xcor() - 20)

screen.onkey(go_left, "Left")   # when Left is pressed, call go_left
screen.listen()                 # switch on the "doorbell" for keys
\end{pycode}
Notice: \texttt{go\_left} has \textbf{no parentheses} in the \texttt{onkey} line. We hand turtle the function, and turtle calls it later.
With parentheses, Python would call it right away, only once.

\subsection{The game loop}
A game is like a flip book: every \textbf{frame} the world moves a tiny bit. Our game loop does one frame, then asks turtle to call it again soon.

\begin{center}
%% alt: The game loop as a cycle of four steps; one lap is one frame, 50 milliseconds. Step 1: move every meteor down. Step 2: check whether anything is touching or off the bottom; if it hit, it is game over and the timer is not set again. Step 3: update the score and the screen. Step 4: set the timer with ontimer(game_loop, 50), which starts the next lap.
\begin{tikzpicture}[font=\small,>={Stealth[length=3mm]},
  st/.style={draw=black,thick,rounded corners=3pt,minimum width=3.6cm,minimum height=0.9cm,align=center}]
  \node[st,fill=kcsky!25] (a) at (0,1.6) {\textbf{1. Move}\\every meteor down};
  \node[st,fill=kcyellow!35] (b) at (4.6,0) {\textbf{2. Check}\\touching? off the bottom?};
  \node[st,fill=kcgreen!25] (c) at (0,-1.6) {\textbf{3. Update}\\score, \texttt{screen.update()}};
  \node[st,fill=kcorange!25] (d) at (-4.6,0) {\textbf{4. Set the timer}\\\texttt{ontimer(game\_loop, 50)}};
  \draw[->,very thick] (a) to[bend left=15] (b);
  \draw[->,very thick] (b) to[bend left=15] (c);
  \draw[->,very thick] (c) to[bend left=15] (d);
  \draw[->,very thick] (d) to[bend left=15] (a);
  \node[draw=kcvermillion,thick,dashed,rounded corners=3pt,align=center,fill=kcvermillion!10] (go) at (4.8,-2.1) {hit? \textbf{GAME OVER}\\(do not set the timer)};
  \draw[->,thick,dashed,kcvermillion] (b) -- (go);
  \node[align=center,font=\footnotesize] at (0,0) {one lap\\= one frame\\(50 ms)};
\end{tikzpicture}
\end{center}

\begin{pycode}[One frame of Dodge the Meteors (shortened)]
state = {"score": 0, "running": True}

def game_loop():
    for meteor in meteors:                      # a list of turtles
        meteor.sety(meteor.ycor() - 8)          # 1. move
        if is_touching(meteor.xcor(), meteor.ycor(),
                       player.xcor(), player.ycor(), 30):
            state["running"] = False            # 2. check: hit!
    if state["running"]:
        screen.ontimer(game_loop, 50)           # 4. again in 50 ms
\end{pycode}

\subsection{Try it: keys and a timer in one small program}
Make a new file in Thonny (\menu{File > New}), type this program, save it as \texttt{try\_keys.py} and press \menu{Run}.
\begin{pycode}[try\_keys.py]
import turtle

screen = turtle.Screen()
player = turtle.Turtle("turtle")
player.penup()
ball = turtle.Turtle("circle")
ball.penup()
ball.goto(0, 200)

def go_left():                    # a handler for the Left key
    player.setx(player.xcor() - 20)

def go_right():                   # a handler for the Right key
    player.setx(player.xcor() + 20)

def fall():                       # one frame: the ball drops a little
    ball.sety(ball.ycor() - 5)
    if ball.ycor() > -200:        # still on the screen?
        screen.ontimer(fall, 50)  # then do it again in 50 ms

screen.onkey(go_left, "Left")
screen.onkey(go_right, "Right")
screen.listen()
fall()
turtle.done()
\end{pycode}
\textbf{What you should see:} the ball falls slowly from the top and stops near the bottom after about 4 seconds.
\textbf{Click the turtle window once}, then press the Left and Right arrows: the turtle moves 20 steps each time.
Nothing happens when you press a key? Click the window, and check \texttt{screen.listen()}.

\begin{tip}[Math helper: up, down, and 50 milliseconds]
\begin{itemize}[leftmargin=1.2em,itemsep=0pt]
  \item The middle of the window is \texttt{(0, 0)}. Going \textbf{down} makes \texttt{y} smaller: 200, 195, 190 \ldots{} 0, then
        below zero: $-5$, $-10$ \ldots{} $-200$ (see the week 1 Math helper). So \texttt{ycor() - 5} moves a thing down.
  \item Going \textbf{left} makes \texttt{x} smaller in the same way: \texttt{xcor() - 20}.
  \item A \textbf{millisecond} (ms) is a tiny part of a second: 1000 milliseconds make 1 second.
        So 50 ms is $\frac{50}{1000}$ of a second, and $1000 \div 50 = 20$ frames every second.
\end{itemize}
\end{tip}

\subsection{Touching = close enough}
Two things touch when the distance between them is small. The \textbf{distance} is how far apart they are in a straight line,
as if you measured from one to the other with a ruler. You do not have to work it out: \texttt{math.dist} measures it for you.
\py|math.dist((0, 0), (30, 40))| gives \texttt{50.0}: the point (30, 40) is 50 steps away from the middle.
If the distance is less than about 30, the meteor hit the player.
Turtles have the same measure built in: \py|player.distance(meteor)| is the distance between two turtles, so
\py|player.distance(meteor) < 30| also means \enquote{touching}.

\begin{exercise}[Etch-a-Sketch]
Folder: \texttt{exercises/week09/ex-etch-a-sketch}. Run \texttt{starter.py}, \textbf{click the window}, press the Up arrow.
Nothing happens yet: the window is not listening.
\begin{enumerate}[leftmargin=1.4em]
  \item TODO 1: add \texttt{screen.listen()}. Now Up draws a line.
  \item TODO 2: write \texttt{go\_down()}, \texttt{go\_left()}, \texttt{go\_right()} (headings 270, 180, 0).
  \item TODO 3: connect the keys: \texttt{"Down"}, \texttt{"Left"}, \texttt{"Right"}, \texttt{"space"}, \texttt{"n"}, \texttt{"c"}.
  \item TODO 4: \texttt{c} clears the drawing. Run \texttt{check\_my\_work.py}: all 5 checks should pass.
        Draw something and describe it here: \answerlines{1}
\end{enumerate}
\end{exercise}

\begin{exercise}[Dodge the Meteors]
Folder: \texttt{exercises/week09/ex-dodge-meteors}. You are the sky-blue turtle; the round orange-red meteors fall.
\begin{enumerate}[leftmargin=1.4em]
  \item TODO 1: \texttt{move\_player()} returns the new x, using \texttt{clamp()} so you stop at the walls.
        The direction is \texttt{1} for right and \texttt{-1} for left: $1 \times 20 = 20$ (20 steps right) and
        $-1 \times 20 = -20$ (20 steps left).
  \item TODO 2: \texttt{is\_touching()} uses \texttt{math.dist} and \texttt{<}. Now run \texttt{check\_my\_work.py}: the first two checks
        (move player, is touching) should say \texttt{[PASS]}.
  \item TODO 3: connect \texttt{"Left"} and \texttt{"Right"}, then \texttt{listen()}.
  \item TODO 4: at the end of \texttt{game\_loop()}, set the timer again. Now play! My best score: \rule{2cm}{0.4pt}
  \item Which functions are \textbf{logic} (no turtle) and which are \textbf{drawing}? Why is it useful to keep them apart? \answerlines{2}
\end{enumerate}
\end{exercise}

\section{Your final project starts now!}
For the next three weeks you will build your own program, alone or with a partner. Read the \textbf{Project Brief} and choose:
\begin{center}
\renewcommand{\arraystretch}{1.3}
\begin{tabular}{@{}p{6.6cm}p{6.6cm}@{}}
\toprule
\textbf{Option A: Arcade Game} & \textbf{Option B: Art Studio} \\ \midrule
A turtle game with keys, a timer and a score, like Dodge the Meteors or Star Catcher. The file saves high scores. &
A program that makes designs (often random) and can save and load them, like Dot Art from a File. \\
\bottomrule
\end{tabular}
\end{center}
Both must use: functions that return values, a list, a dictionary, a loop or game loop, a file, and at least 3 tests.
\textbf{Milestone M1:} Planning page 1 (your pitch) is homework this week. Page 2 (pseudocode or flowchart) and, in a pair,
page 3 (task split) are finished together at the start of next week; you may start page 2 at home if you like.

\begin{tip}[Borrowing pictures or sounds]
Pictures and sounds belong to the people who made them. Only use ones with an open license (CC0: free for anyone; CC BY: free if you give credit),
ask a grown-up to help you find them, and write a credit in your code, like \texttt{\# Star sound by Mei (CC BY)}. Turtle's own shapes are always fine!
\end{tip}

\begin{vocab}
  \vocabword{event}{Something that happens while a program runs: a key press, a click, a timer ringing.}
  \vocabword{handler}{The function that runs when an event happens.}
  \vocabword{listen()}{Tells the turtle window to start hearing key presses.}
  \vocabword{game loop}{A function that does one frame of the game, then sets a timer to run again.}
  \vocabword{frame}{One tiny step of a game: everything moves a little, then it is drawn.}
  \vocabword{collision}{Two things touching. We check it with the distance between them.}
  \vocabword{license / attribution}{The rules for using someone's work / giving them credit for it.}
\end{vocab}

\begin{reflection}
What happens if you forget \texttt{ontimer} at the end of the game loop? \answerlines{1}
Explain to a friend why \texttt{onkey(go\_left(), "Left")} does not work. \answerlines{2}
My project idea in one sentence: \answerlines{1}
\end{reflection}
